% Main thesis file
\documentclass[12pt,a4paper,twoside,openright]{scrbook}

% ============================================
% PACKAGES
% ============================================

% Language and encoding
\usepackage[utf8]{inputenc}
\usepackage[T1]{fontenc}
\usepackage[english]{babel}

% Bibliography
\usepackage[backend=biber,style=authoryear-ibid,sorting=nyt,maxbibnames=99]{biblatex}
\addbibresource{references.bib}

% Math packages
\usepackage{amsmath,amssymb,amsthm}

% Graphics and figures
\usepackage{graphicx}
\usepackage{float}
\usepackage[font=small,labelfont=bf]{caption}
\usepackage{subcaption}

% Tables
\usepackage{booktabs}
\usepackage{multirow}
\usepackage{array}

% Colors and boxes
\usepackage{xcolor}
\usepackage{tcolorbox}

% Hyperlinks (load last)
\usepackage[hidelinks]{hyperref}
\usepackage{cleveref}

% Code listings
\usepackage{listings}
\lstset{
    basicstyle=\ttfamily\small,
    breaklines=true,
    frame=single,
    numbers=left,
    numberstyle=\tiny,
    captionpos=b
}

% Acronyms
\usepackage[acronym,toc]{glossaries}
\makeglossaries

% Define acronyms
\newacronym{ssl}{SSL}{Self-Supervised Learning}
\newacronym{cnn}{CNN}{Convolutional Neural Network}
\newacronym{hil}{HIL}{Hardware-in-the-Loop}
\newacronym{adas}{ADAS}{Advanced Driver Assistance Systems}
\newacronym{can}{CAN}{Controller Area Network}
\newacronym{obd}{OBD}{On-Board Diagnostics}
\newacronym{ml}{ML}{Machine Learning}
\newacronym{ai}{AI}{Artificial Intelligence}
\newacronym{auc}{AUC}{Area Under the Curve}
\newacronym{roc}{ROC}{Receiver Operating Characteristic}

% ============================================
% DOCUMENT INFORMATION
% ============================================

\title{Automotive Sensors Fault Detection and Localization Using Self-Supervised Learning}
\author{Your Name}
\date{\today}

% University information (customize as needed)
\newcommand{\university}{Your University Name}
\newcommand{\department}{Department of Computer Science}
\newcommand{\degreetype}{Master of Science}
\newcommand{\supervisor}{Prof. Supervisor Name}

% ============================================
% DOCUMENT START
% ============================================

\begin{document}

% ============================================
% FRONT MATTER
% ============================================

\frontmatter

% Title page
\begin{titlepage}
    \centering
    \vspace*{2cm}

    {\LARGE \university \par}
    \vspace{0.5cm}
    {\Large \department \par}

    \vspace{3cm}

    {\huge\bfseries Automotive Sensors Fault Detection and Localization Using Self-Supervised Learning\par}

    \vspace{2cm}

    {\Large A thesis submitted in partial fulfillment of the requirements\\
    for the degree of \degreetype\par}

    \vspace{2cm}

    {\large
    \textbf{Author:}\\
    Your Name\\[0.5cm]

    \textbf{Supervisor:}\\
    \supervisor\\[0.5cm]

    \textbf{Date:}\\
    \today
    }

    \vfill
\end{titlepage}

% Declaration of authorship
\chapter*{Declaration of Authorship}
\addcontentsline{toc}{chapter}{Declaration of Authorship}

I hereby declare that this thesis is my own work and that I have not used any sources other than those listed in the bibliography and identified as references. I further declare that I have not submitted this thesis at any other institution in order to obtain a degree.

\vspace{2cm}

\noindent
\rule{5cm}{0.4pt} \hfill \rule{5cm}{0.4pt}\\
Date \hfill Signature

\cleardoublepage

% Abstract (English)
\chapter*{Abstract}
\addcontentsline{toc}{chapter}{Abstract}

Modern vehicles rely on numerous sensors for \gls{adas} and safety-critical functions. Sensor faults can lead to system failures, compromising vehicle safety and performance. This thesis presents an adaptive \gls{ssl}-based approach for detecting and localizing sensor faults in automotive systems without requiring labeled fault data during training.

We propose a framework that employs a 1D \gls{cnn} encoder trained on transformation classification tasks using normal sensor data. The learned representations are used for fault detection via Mahalanobis distance-based anomaly detection and fault localization through ablation-based analysis. A key contribution is the adaptive learning capability, where the system incrementally updates its statistical models and fault classifier without retraining the encoder, enabling continuous evolution with new operational data.

The approach was evaluated on the Audi A2D2 dataset with \gls{hil}-calibrated fault injection across multiple sensor types. Results demonstrate 77.5\% detection accuracy, 67.6\% top-1 localization accuracy (85.3\% top-3), and effective adaptation to new fault patterns. The framework requires no labeled fault examples during initial training, making it practical for real-world automotive applications where fault data is scarce.

\textbf{Keywords:} Self-supervised learning, Fault detection, Automotive sensors, Anomaly detection, Adaptive learning, Hardware-in-the-loop

\cleardoublepage

% Table of contents
\tableofcontents
\cleardoublepage

% List of figures
\listoffigures
\addcontentsline{toc}{chapter}{List of Figures}
\cleardoublepage

% List of tables
\listoftables
\addcontentsline{toc}{chapter}{List of Tables}
\cleardoublepage

% List of acronyms
\printglossary[type=\acronymtype,title=List of Acronyms]
\cleardoublepage

% ============================================
% MAIN MATTER - CHAPTERS
% ============================================

\mainmatter

% Include chapter files
\chapter{Introduction}
\label{chap:introduction}

\section{Motivation}

Modern vehicles have evolved into complex cyber-physical systems, incorporating hundreds of sensors to enable \gls{adas}, improve safety, and enhance user experience. These sensors monitor critical parameters such as speed, acceleration, steering angle, brake pressure, and environmental conditions. The reliability of these sensors is paramount, as sensor faults can lead to incorrect decision-making by automated systems, potentially compromising vehicle safety and performance \parencite{iso26262_2018}.

According to the ISO 26262 functional safety standard, automotive systems must detect and respond to sensor failures to maintain safe operation \parencite{iso26262_2018}. Traditional fault detection approaches rely on physical redundancy (duplicate sensors) or analytical redundancy (mathematical models), but these methods have limitations: physical redundancy increases cost and complexity, while model-based approaches require extensive domain knowledge and may not generalize well to diverse operating conditions.

Recent advances in \gls{ml} and \gls{ai} offer promising alternatives for sensor fault detection. However, supervised learning approaches require large labeled datasets of fault conditions, which are difficult and expensive to collect in real-world automotive applications. Intentionally inducing sensor faults in operational vehicles poses safety risks, and naturally occurring faults are rare events that may not cover all failure modes.

\section{Problem Statement}

The key challenge addressed in this thesis is: \textit{How can we detect and localize sensor faults in automotive systems without requiring extensive labeled fault data, while enabling the system to adapt continuously to new operational patterns?}

This problem is particularly relevant because:

\begin{itemize}
    \item \textbf{Data scarcity}: Labeled fault data is expensive to acquire and may not represent all possible fault scenarios
    \item \textbf{Operational variability}: Vehicles operate under diverse conditions (weather, driving styles, road types) that affect sensor behavior
    \item \textbf{Evolving systems}: Automotive systems are updated over their lifetime, requiring fault detection to adapt without complete retraining
    \item \textbf{Real-time constraints}: Fault detection must operate efficiently on embedded automotive hardware
\end{itemize}

\section{Research Objectives}

The primary objective of this thesis is to develop an adaptive fault detection and localization framework for automotive sensors using \gls{ssl}. Specific objectives include:

\begin{enumerate}
    \item \textbf{Self-supervised representation learning}: Develop a method to learn meaningful sensor data representations using only normal (fault-free) operational data
    \item \textbf{Fault detection without labeled examples}: Design an anomaly detection mechanism that can identify sensor faults without requiring fault examples during training
    \item \textbf{Fault localization}: Implement a technique to identify which specific sensor(s) are faulty when an anomaly is detected
    \item \textbf{Adaptive learning capability}: Enable the system to update its models incrementally as new operational data becomes available, without retraining from scratch
    \item \textbf{Empirical validation}: Evaluate the approach on real automotive sensor data with realistic fault injection scenarios
\end{enumerate}

\section{Contributions}

This thesis makes the following contributions:

\begin{itemize}
    \item \textbf{SSL-based automotive fault detection framework}: A novel application of self-supervised learning to automotive sensor fault detection, using transformation classification as the pretext task

    \item \textbf{Adaptive anomaly detection}: An incremental learning mechanism for Mahalanobis distance-based anomaly detection that updates statistical models without retraining the feature encoder

    \item \textbf{Incremental fault classifier}: A framework for continuously learning new fault patterns using online learning algorithms, enabling adaptation to previously unseen fault types

    \item \textbf{Ablation-based fault localization}: A method to identify faulty sensors by systematically masking individual sensors and analyzing the impact on anomaly scores

    \item \textbf{HIL-calibrated evaluation}: Comprehensive evaluation using Hardware-in-the-Loop calibrated fault ranges on the Audi A2D2 dataset, demonstrating practical applicability

    \item \textbf{Statistical validation}: Rigorous statistical analysis including confidence intervals, precision-recall metrics, and ROC-AUC analysis across multiple experimental runs
\end{itemize}

\section{Thesis Structure}

The remainder of this thesis is organized as follows:

\begin{description}
    \item[Chapter \ref{chap:related_work}] reviews related work on self-supervised learning, anomaly detection, automotive fault diagnosis, and incremental learning approaches.

    \item[Chapter \ref{chap:methodology}] describes the proposed methodology, including the SSL pretext task, the 1D CNN encoder architecture, Mahalanobis distance-based anomaly detection, ablation-based localization, and the adaptive learning framework.

    \item[Chapter \ref{chap:implementation}] details the implementation, including dataset preparation, model architecture specifics, training procedures, and fault injection methodology.

    \item[Chapter \ref{chap:results}] presents experimental results, including detection accuracy, localization performance, ablation studies, and comparison with baselines.

    \item[Chapter \ref{chap:conclusion}] concludes the thesis with a summary of findings, limitations of the current approach, and directions for future research.
\end{description}

\chapter{Related Work}
\label{chap:related_work}

This chapter reviews prior research relevant to our approach, organized into four main areas: self-supervised learning, anomaly detection for time-series data, automotive fault diagnosis, and incremental learning methods.

\section{Self-Supervised Learning}

\gls{ssl} has emerged as a powerful paradigm for learning useful representations from unlabeled data \parencite{lecun2015deep}. Unlike supervised learning, which requires labeled examples, SSL formulates pretext tasks that generate supervisory signals from the data itself.

\subsection{Contrastive Learning Approaches}

Contrastive learning methods have achieved remarkable success in computer vision. SimCLR \parencite{chen2020simple} learns representations by maximizing agreement between differently augmented views of the same image while minimizing agreement with other images. MoCo \parencite{he2020momentum} uses a momentum encoder and queue mechanism to maintain a large and consistent dictionary for contrastive learning. These approaches have demonstrated that representations learned without labels can match or exceed supervised learning performance on downstream tasks.

While highly successful in vision domains, direct application of contrastive learning to time-series sensor data faces challenges due to the temporal nature of the data and the difficulty of defining meaningful augmentations that preserve semantic content.

\subsection{Prediction-Based Methods}

Another SSL approach involves predicting missing or future parts of the input. Contrastive Predictive Coding (CPC) \parencite{oord2018representation} learns representations by predicting future observations in latent space using contrastive objectives. BYOL \parencite{grill2020bootstrap} eliminates the need for negative samples by using a momentum encoder and predictor network.

For automotive sensor data, prediction-based methods can leverage the temporal correlations inherent in vehicle dynamics, but may struggle with the high variability of driving scenarios.

\subsection{Transformation-Based Pretext Tasks}

An alternative SSL strategy is to train models to recognize transformations applied to the data. For time-series, \textcite{tran2022self} propose using transformations such as jittering, masking, and scaling as pretext tasks for anomaly detection in industrial IoT systems. The intuition is that a model trained to recognize these transformations on normal data will produce unusual representations for anomalous data.

Our work builds on this transformation-based approach, adapting it specifically for automotive sensor data with multiple channels and developing an adaptive framework that can evolve over time.

\section{Anomaly Detection for Time-Series}

Anomaly detection in time-series data is a well-studied problem with applications across many domains \parencite{salehi2021unified}.

\subsection{Statistical Methods}

Traditional statistical approaches model normal data distribution and flag observations that deviate significantly. Methods include z-score analysis, isolation forests, and one-class SVM. While interpretable, these methods often struggle with high-dimensional multivariate time-series and complex normal variability.

Mahalanobis distance-based approaches \parencite{de2000mahalanobis} account for correlations between variables and have been successfully applied to out-of-distribution detection in neural network representations \parencite{lee2018simple}. We adopt this approach for its theoretical foundation and computational efficiency.

\subsection{Deep Learning Methods}

Deep learning has enabled more sophisticated anomaly detection. Autoencoders learn compressed representations and detect anomalies via reconstruction error. \textcite{su2019robust} propose a stochastic recurrent neural network for multivariate time-series anomaly detection. USAD \parencite{audibert2020usad} uses adversarial training with autoencoders for unsupervised anomaly detection.

LSTM-based approaches can model temporal dependencies but require careful architecture design and may be computationally expensive for real-time automotive applications. Our CNN-based approach offers a favorable trade-off between modeling capacity and computational efficiency.

\section{Automotive Fault Diagnosis}

Fault detection and diagnosis in automotive systems has attracted significant research attention due to safety-critical implications.

\subsection{Model-Based Approaches}

Traditional automotive fault detection relies on analytical redundancy, where mathematical models of expected sensor behavior are compared with actual observations \parencite{iso26262_2018}. While these methods are well-understood and verifiable, they require extensive domain expertise and may not generalize to all operating conditions.

\subsection{Data-Driven Methods}

Machine learning approaches have shown promise for automotive fault diagnosis. \textcite{zhao2019machine} provide a comprehensive survey of deep learning for machine health monitoring. \textcite{wang2020deep} review deep learning-based fault diagnosis across various domains.

Most existing work focuses on supervised learning with labeled fault data. \textcite{lu2020intelligent} use hierarchical CNNs for bearing fault classification, achieving high accuracy but requiring extensive labeled fault examples. The challenge of obtaining labeled fault data in automotive contexts motivates our self-supervised approach.

\subsection{Hardware-in-the-Loop Testing}

\gls{hil} simulation is widely used in automotive development for testing control systems under realistic conditions without requiring a physical vehicle \parencite{fathy2006review, bringmann2008model}. HIL systems can inject faults in a controlled manner, providing valuable data for fault detection algorithm development. We leverage HIL-calibrated fault ranges to create realistic test scenarios.

\section{Incremental and Online Learning}

Adaptive systems that can learn continuously without forgetting are crucial for long-term autonomous operation.

\subsection{Incremental Learning Algorithms}

Incremental learning methods update models with new data without retraining from scratch. \textcite{losing2018incremental} provide a comprehensive review of online learning algorithms. Stochastic Gradient Descent (SGD) based approaches \parencite{pmlr2011online} enable efficient updates with streaming data.

Our incremental fault classifier uses SGD with warm start, allowing new fault types to be learned while retaining knowledge of previously seen faults.

\subsection{Incremental Statistics}

For anomaly detection, incrementally updating distributional statistics is essential. Welford's algorithm \parencite{welford1962note} provides a numerically stable method for computing mean and variance incrementally. We extend this to covariance matrices for multivariate Mahalanobis distance computation.

\section{Datasets}

The Audi A2D2 dataset \parencite{geiger2020a2d2} provides real-world automotive sensor data collected from an Audi vehicle equipped with various sensors including cameras, LiDAR, and vehicle bus signals. The dataset includes diverse driving scenarios and is well-suited for developing and evaluating automotive perception and fault detection algorithms. We use the vehicle bus signal data, which contains \gls{can} bus sensor readings at high frequency.

\section{Summary}

While existing research has made significant progress in SSL, anomaly detection, and automotive fault diagnosis individually, the combination of these techniques into an adaptive framework specifically designed for automotive sensor fault detection represents a novel contribution. Our work bridges these areas by:

\begin{itemize}
    \item Applying transformation-based SSL to automotive multivariate sensor time-series
    \item Using Mahalanobis distance on learned representations for anomaly detection
    \item Developing an incremental learning framework that adapts without encoder retraining
    \item Demonstrating practical applicability with HIL-calibrated fault injection
\end{itemize}

The following chapter describes our methodology in detail.

\chapter{Methodology}
\label{chap:methodology}

This chapter presents our adaptive self-supervised learning framework for automotive sensor fault detection and localization. We first provide an overview of the system architecture, then detail each component.

\section{Framework Overview}

\Cref{fig:framework_overview} illustrates the overall framework, which consists of four main stages:

\begin{enumerate}
    \item \textbf{Self-Supervised Representation Learning}: A 1D \gls{cnn} encoder is trained on normal sensor data using transformation classification as the pretext task
    \item \textbf{Feature Extraction and Statistical Modeling}: The trained encoder extracts features from normal data, and distributional statistics (mean, covariance) are computed for Mahalanobis distance calculation
    \item \textbf{Fault Detection}: New sensor data is encoded and its Mahalanobis distance from the normal distribution is computed; distances exceeding a threshold indicate potential faults
    \item \textbf{Fault Localization}: When a fault is detected, ablation-based analysis identifies which sensor(s) are faulty
\end{enumerate}

A key innovation is the \textit{adaptive learning capability}: the system can incrementally update its statistical models and fault classifier with new data without retraining the encoder, enabling continuous evolution in operational environments.

% Note: Add figure in final version
% \begin{figure}[htbp]
%     \centering
%     \includegraphics[width=0.9\textwidth]{figures/framework_overview.pdf}
%     \caption{Overview of the adaptive SSL-based fault detection framework}
%     \label{fig:framework_overview}
% \end{figure}

\section{Self-Supervised Representation Learning}

\subsection{Motivation for Self-Supervised Learning}

Supervised fault detection requires labeled examples of various fault conditions, which are difficult to obtain in automotive contexts. Self-supervised learning addresses this by learning useful representations from normal operational data alone. The hypothesis is that representations learned to recognize transformations of normal data will produce distinctive patterns for abnormal (faulty) data.

\subsection{Transformation Classification Pretext Task}

We define three transformations applied to sensor time-series windows:

\begin{description}
    \item[Jitter (T1)] Adds small Gaussian noise: $\mathbf{x}' = \mathbf{x} + \mathcal{N}(0, \sigma^2)$

    \item[Masking (T2)] Randomly sets segments to zero: $x'_i = 0$ for $i \in \mathcal{M}$, where $\mathcal{M}$ is a random subset of time steps

    \item[Negation (T3)] Inverts the signal: $\mathbf{x}' = -\mathbf{x}$
\end{description}

These transformations are chosen because:
\begin{itemize}
    \item They preserve semantic content while creating distinguishable variations
    \item They are relevant to automotive scenarios (jitter simulates sensor noise, masking simulates signal dropout, negation captures polarity)
    \item They enable a simple multi-class classification task suitable for CNNs
\end{itemize}

The pretext task is to classify which transformation was applied to a given window. Let $\mathbf{x} \in \mathbb{R}^{T \times C}$ be a time-series window with $T$ time steps and $C$ channels (sensors). The encoder $f_\theta: \mathbb{R}^{T \times C} \rightarrow \mathbb{R}^{d}$ maps input to a $d$-dimensional representation. A classifier head $g_\phi: \mathbb{R}^{d} \rightarrow \mathbb{R}^{K}$ predicts the transformation class ($K=3$ transformations). Training minimizes cross-entropy loss:

\begin{equation}
\mathcal{L}_{\text{SSL}} = -\sum_{i=1}^{N} \sum_{k=1}^{K} y_{ik} \log(\hat{y}_{ik})
\end{equation}

where $y_{ik}$ is the true label (one-hot encoded) and $\hat{y}_{ik} = \text{softmax}(g_\phi(f_\theta(\mathbf{x}_i)))_k$.

\subsection{Encoder Architecture}

We use a 1D \gls{cnn} architecture \parencite{kiranyaz20191d} suited for time-series data. The encoder consists of:

\begin{itemize}
    \item \textbf{Input layer}: Accepts windows of size $(T, C)$
    \item \textbf{Convolutional layers}: Three 1D convolutional layers with kernel sizes (7, 5, 3), stride 1, and increasing channels (64, 128, 256)
    \item \textbf{Activation}: ReLU activation after each convolutional layer
    \item \textbf{Batch normalization}: Applied after each activation for training stability
    \item \textbf{Global average pooling}: Aggregates temporal dimension to fixed-size representation
    \item \textbf{Output}: $d$-dimensional embedding (typically $d=256$)
\end{itemize}

The relatively shallow architecture balances representational capacity with computational efficiency, which is crucial for real-time operation on embedded automotive hardware.

\section{Fault Detection via Mahalanobis Distance}

\subsection{Distributional Modeling}

After training the encoder on normal data, we extract embeddings $\{\mathbf{z}_i\}_{i=1}^{N}$ from a validation set of normal samples, where $\mathbf{z}_i = f_\theta(\mathbf{x}_i) \in \mathbb{R}^{d}$. We compute the sample mean and covariance:

\begin{align}
\boldsymbol{\mu} &= \frac{1}{N} \sum_{i=1}^{N} \mathbf{z}_i \\
\boldsymbol{\Sigma} &= \frac{1}{N-1} \sum_{i=1}^{N} (\mathbf{z}_i - \boldsymbol{\mu})(\mathbf{z}_i - \boldsymbol{\mu})^\top
\end{align}

\subsection{Mahalanobis Distance for Anomaly Detection}

The Mahalanobis distance \parencite{de2000mahalanobis} measures how far a point $\mathbf{z}$ is from the distribution center, accounting for correlations:

\begin{equation}
D_M(\mathbf{z}) = \sqrt{(\mathbf{z} - \boldsymbol{\mu})^\top \boldsymbol{\Sigma}^{-1} (\mathbf{z} - \boldsymbol{\mu})}
\end{equation}

For numerical stability, we use the squared Mahalanobis distance $D_M^2(\mathbf{z})$ and add a regularization term $\epsilon \mathbf{I}$ to the covariance matrix:

\begin{equation}
\boldsymbol{\Sigma}_{\text{reg}} = \boldsymbol{\Sigma} + \epsilon \mathbf{I}
\end{equation}

where $\epsilon = 10^{-6}$ prevents singularity.

A sample is classified as anomalous (faulty) if:

\begin{equation}
D_M^2(\mathbf{z}) > \tau
\end{equation}

where $\tau$ is a threshold determined from the validation set. We set $\tau$ to the 90th percentile of validation Mahalanobis distances, ensuring a 10\% false positive rate on normal data.

\subsection{Rationale for Mahalanobis Distance}

We chose Mahalanobis distance over simpler metrics (Euclidean distance, reconstruction error) because:
\begin{itemize}
    \item It accounts for correlations between embedding dimensions
    \item It has a statistical interpretation (related to the $\chi^2$ distribution under Gaussian assumptions)
    \item It has been shown effective for out-of-distribution detection in neural networks \parencite{lee2018simple}
    \item It is computationally efficient (matrix multiplication) for real-time inference
\end{itemize}

\section{Fault Localization}

Detecting that a fault exists is insufficient; identifying the faulty sensor(s) is critical for effective diagnosis and mitigation.

\subsection{Ablation-Based Localization}

Our localization approach is based on the hypothesis: \textit{removing a faulty sensor from the input should reduce the anomaly score, while removing a healthy sensor should not significantly affect it}.

Given a detected anomaly with embedding $\mathbf{z}_{\text{full}}$ (all sensors), we systematically create ablated versions by masking each sensor individually:

\begin{equation}
\mathbf{x}^{(-c)} = \mathbf{x} \odot \mathbf{m}^{(-c)}
\end{equation}

where $\mathbf{m}^{(-c)}$ is a mask that zeros out sensor $c$, and $\odot$ is element-wise multiplication. We compute the Mahalanobis distance for each ablated input:

\begin{equation}
D_c = D_M^2(f_\theta(\mathbf{x}^{(-c)}))
\end{equation}

The localization score for sensor $c$ is the reduction in anomaly score when that sensor is removed:

\begin{equation}
S_c = D_M^2(\mathbf{z}_{\text{full}}) - D_c
\end{equation}

Higher $S_c$ indicates sensor $c$ is more likely faulty. We rank sensors by $S_c$ and report the top-1 (most likely faulty) and top-3 sensors.

\subsection{Multi-Sensor Fault Localization}

In cases where multiple sensors are faulty, we can extend the approach by considering combinations of sensors. However, this increases computational cost exponentially with the number of sensors. For efficiency, we use the single-sensor ablation approach and leave multi-sensor combinatorial analysis for future work.

\section{Adaptive Learning Framework}

A key contribution of our work is enabling the system to adapt over time without full retraining, which is critical for long-term deployment in operational vehicles.

\subsection{Frozen Encoder, Updateable Statistics}

Once the encoder $f_\theta$ is trained, we freeze its parameters. Adaptation occurs by updating the distributional statistics ($\boldsymbol{\mu}$, $\boldsymbol{\Sigma}$) and the fault classifier, while keeping the encoder fixed. This design choice is motivated by:

\begin{itemize}
    \item \textbf{Stability}: The encoder learned general representations from diverse normal data; updating it with limited new data risks overfitting or catastrophic forgetting
    \item \textbf{Efficiency}: Updating statistics is computationally cheap compared to retraining a neural network
    \item \textbf{Interpretability}: Changes in anomaly detection behavior are attributable to distributional shifts, not representation changes
\end{itemize}

\subsection{Incremental Statistics Update}

When new normal data $\{\mathbf{z}_{\text{new},j}\}_{j=1}^{M}$ becomes available, we update the mean and covariance incrementally using Welford's algorithm \parencite{welford1962note}.

Let $\boldsymbol{\mu}_{\text{old}}$, $\boldsymbol{\Sigma}_{\text{old}}$, and $n_{\text{old}}$ be the previous mean, covariance, and sample count. The updated statistics are:

\begin{align}
n_{\text{new}} &= n_{\text{old}} + M \\
\boldsymbol{\mu}_{\text{new}} &= \frac{n_{\text{old}} \boldsymbol{\mu}_{\text{old}} + M \bar{\mathbf{z}}_{\text{new}}}{n_{\text{new}}}
\end{align}

where $\bar{\mathbf{z}}_{\text{new}} = \frac{1}{M} \sum_{j=1}^{M} \mathbf{z}_{\text{new},j}$ is the mean of new embeddings.

For covariance, we use the parallel algorithm:

\begin{equation}
\boldsymbol{\Sigma}_{\text{new}} = \frac{(n_{\text{old}}-1)\boldsymbol{\Sigma}_{\text{old}} + (M-1)\boldsymbol{\Sigma}_{\text{new}}' + \frac{n_{\text{old}} M}{n_{\text{new}}}(\boldsymbol{\mu}_{\text{old}} - \bar{\mathbf{z}}_{\text{new}})(\boldsymbol{\mu}_{\text{old}} - \bar{\mathbf{z}}_{\text{new}})^\top}{n_{\text{new}}-1}
\end{equation}

where $\boldsymbol{\Sigma}_{\text{new}}'$ is the covariance of the new samples.

This update is numerically stable and computationally efficient ($O(d^2)$ operations).

\subsection{Incremental Fault Classifier}

In addition to anomaly detection, we maintain a fault classifier that learns to categorize detected faults by type (e.g., offset, noise, dropout). This classifier uses Stochastic Gradient Descent with warm start \parencite{pmlr2011online}, enabling incremental learning.

When new labeled fault data becomes available (e.g., from service records or HIL testing), we perform a partial fit:

\begin{equation}
\theta_{\text{classifier}} \leftarrow \theta_{\text{classifier}} - \eta \nabla_\theta \mathcal{L}(\mathbf{z}, y)
\end{equation}

where $\eta$ is the learning rate, $\mathbf{z}$ is the embedding of a faulty sample, $y$ is its fault type label, and $\mathcal{L}$ is the log loss.

The classifier can learn new fault types dynamically by specifying the full class set during the first update, then adding samples from new classes as they become available.

\subsection{Adaptation Workflow}

The operational adaptation workflow is:

\begin{enumerate}
    \item \textbf{Continuous monitoring}: System processes incoming sensor data and performs fault detection
    \item \textbf{Normal data collection}: Samples classified as normal with high confidence are stored
    \item \textbf{Periodic statistics update}: At regular intervals (e.g., daily), accumulated normal samples are used to update $\boldsymbol{\mu}$ and $\boldsymbol{\Sigma}$
    \item \textbf{Fault data labeling}: Detected faults are logged; a subset may be labeled (e.g., through service technician feedback or HIL validation)
    \item \textbf{Classifier update}: Labeled fault samples are used to update the fault classifier incrementally
\end{enumerate}

This workflow enables the system to adapt to gradual distribution shifts (e.g., component aging, environmental changes) and learn new fault patterns without human intervention for normal operation.

\section{Ensemble Methods for Robustness}

To improve robustness, we train multiple encoder models with different random initializations and aggregate their predictions. For detection, we average the Mahalanobis distances:

\begin{equation}
D_{\text{ensemble}}^2(\mathbf{z}) = \frac{1}{K} \sum_{k=1}^{K} D_{M,k}^2(\mathbf{z})
\end{equation}

For localization, we average the localization scores across models. Ensemble methods reduce variance and improve generalization.

\section{Summary}

Our methodology combines:
\begin{itemize}
    \item \textbf{Self-supervised learning} to learn representations without labeled fault data
    \item \textbf{Mahalanobis distance} for statistically-grounded anomaly detection
    \item \textbf{Ablation-based localization} to identify faulty sensors
    \item \textbf{Incremental adaptation} to evolve with new data without full retraining
\end{itemize}

The following chapter describes the implementation details and experimental setup.

\chapter{Implementation}
\label{chap:implementation}

This chapter describes the implementation details of our adaptive fault detection framework, including dataset preparation, model architecture specifics, training procedures, and fault injection methodology.

\section{Dataset}

\subsection{Audi A2D2 Dataset}

We use the Audi Autonomous Driving Dataset (A2D2) \parencite{geiger2020a2d2}, which contains sensor data from an instrumented Audi test vehicle. The dataset includes:

\begin{itemize}
    \item Vehicle bus signals sampled at 100 Hz
    \item Diverse driving scenarios including urban, highway, and rural environments
    \item Various weather and lighting conditions
    \item Multiple sensor types: speed, acceleration, steering angle, brake pressure, etc.
\end{itemize}

We focus on the vehicle bus signal data, which provides multivariate time-series of sensor readings from the \gls{can} bus.

\subsection{Sensor Selection}

From the available bus signals, we select 8 sensors that are critical for vehicle dynamics and safety:

\begin{enumerate}
    \item \textbf{Longitudinal acceleration} -- acceleration in forward/backward direction
    \item \textbf{Lateral acceleration} -- acceleration in left/right direction
    \item \textbf{Vehicle speed} -- current velocity of the vehicle
    \item \textbf{Steering angle} -- angle of the steering wheel
    \item \textbf{Yaw rate} -- rotational velocity around vertical axis
    \item \textbf{Brake pressure} -- hydraulic pressure in brake system
    \item \textbf{Throttle position} -- accelerator pedal position
    \item \textbf{Wheel speeds} -- rotational speeds of individual wheels
\end{enumerate}

These sensors are chosen because they are commonly available in modern vehicles, are used by ADAS features, and their faults can significantly impact vehicle safety and performance.

\subsection{Data Preprocessing}

Raw sensor data undergoes several preprocessing steps:

\begin{enumerate}
    \item \textbf{Missing value handling}: Interpolate or remove sequences with excessive missing values
    \item \textbf{Normalization}: Apply z-score normalization per sensor channel:
    \begin{equation}
    \tilde{x}_i = \frac{x_i - \mu_i}{\sigma_i}
    \end{equation}
    where $\mu_i$ and $\sigma_i$ are the mean and standard deviation of sensor $i$ computed from training data
    \item \textbf{Windowing}: Segment time-series into fixed-length windows of $T=50$ time steps (0.5 seconds at 100 Hz)
    \item \textbf{Stride}: Use stride of 25 for training (50\% overlap) to increase training samples; use stride of 50 for testing (no overlap) to evaluate independent samples
\end{enumerate}

\subsection{Train-Validation-Test Split}

We split the dataset temporally to simulate realistic deployment:

\begin{itemize}
    \item \textbf{Training set}: First 70\% of driving sessions, used for SSL encoder training
    \item \textbf{Validation set}: Next 15\%, used for threshold calibration and hyperparameter tuning
    \item \textbf{Test set}: Final 15\%, used for final evaluation with fault injection
\end{itemize}

Importantly, only \textit{normal} (fault-free) data is used for training and validation. Faults are injected only into the test set.

\section{Model Architecture and Training}

\subsection{Encoder Architecture Details}

The 1D CNN encoder is implemented in PyTorch with the following architecture:

\begin{table}[htbp]
\centering
\caption{Encoder architecture specifications}
\label{tab:encoder_arch}
\begin{tabular}{@{}lcccc@{}}
\toprule
Layer & Kernel Size & Stride & Channels Out & Activation \\
\midrule
Conv1D-1 & 7 & 1 & 64 & ReLU \\
BatchNorm1D & - & - & 64 & - \\
Conv1D-2 & 5 & 1 & 128 & ReLU \\
BatchNorm1D & - & - & 128 & - \\
Conv1D-3 & 3 & 1 & 256 & ReLU \\
BatchNorm1D & - & - & 256 & - \\
GlobalAvgPool1D & - & - & 256 & - \\
\bottomrule
\end{tabular}
\end{table}

The output is a 256-dimensional embedding vector. A classification head (linear layer + softmax) is attached during SSL training but discarded afterward.

\subsection{Training Procedure}

\paragraph{Pretext Task Training}

The encoder is trained on the transformation classification task with the following hyperparameters:

\begin{itemize}
    \item \textbf{Optimizer}: Adam with learning rate $\eta = 0.001$
    \item \textbf{Batch size}: 128
    \item \textbf{Epochs}: 50 (with early stopping based on validation loss)
    \item \textbf{Loss function}: Cross-entropy
    \item \textbf{Regularization}: Weight decay of $10^{-5}$
    \item \textbf{Transformation parameters}: Jitter noise std $\sigma = 0.1$, masking ratio $p = 0.15$
\end{itemize}

We train 5 models with different random seeds to create an ensemble.

\paragraph{Threshold Calibration}

After training, we extract embeddings from the validation set and compute Mahalanobis distances. The threshold $\tau$ is set to the 90th percentile of these distances, targeting a 10\% false positive rate on normal data.

\subsection{Computational Efficiency}

The encoder has approximately 180K parameters. On a standard GPU (NVIDIA RTX 2060), inference time per window is $\sim$0.5 ms, enabling real-time operation. Training takes approximately 2 hours for 50 epochs.

\section{Fault Injection Methodology}

To evaluate the framework, we inject synthetic faults into the test set. The fault types and parameters are calibrated using \gls{hil} testing data to ensure realism.

\subsection{HIL-Calibrated Fault Ranges}

HIL test data provides realistic bounds for sensor fault magnitudes. We analyze historical HIL logs to determine typical fault ranges:

\begin{table}[htbp]
\centering
\caption{HIL-calibrated fault injection parameters}
\label{tab:fault_params}
\begin{tabular}{@{}lll@{}}
\toprule
Fault Type & Description & Parameter Range \\
\midrule
Offset & Constant bias added to sensor & $[-2\sigma, 2\sigma]$ \\
Scaling & Multiplicative factor & $[0.5, 2.0]$ \\
Noise & Added Gaussian noise & std $\in [0.5\sigma, 2\sigma]$ \\
Drift & Linearly increasing bias & slope $\in [-0.1\sigma, 0.1\sigma]$ per step \\
Dropout & Intermittent signal loss & dropout rate $\in [0.1, 0.5]$ \\
\bottomrule
\end{tabular}
\end{table}

where $\sigma$ is the standard deviation of the sensor's normal values.

\subsection{Multi-Run Fault Injection Protocol}

To ensure statistical validity, we perform multiple fault injection runs:

\begin{enumerate}
    \item \textbf{Random selection}: For each run, randomly select $N_f$ windows from the test set
    \item \textbf{Random fault type}: Assign a random fault type from Table~\ref{tab:fault_params}
    \item \textbf{Random sensor}: Select a random sensor to be faulty (ground truth for localization)
    \item \textbf{Random parameters}: Sample fault parameters uniformly from the calibrated ranges
    \item \textbf{Fault application}: Apply the fault to the selected sensor in the selected window
\end{enumerate}

We perform 15 independent runs with different random seeds, generating a total of 58 faulty samples for robust evaluation.

\subsection{Rationale for Synthetic Faults}

While synthetic faults are not perfect representations of real-world failures, they offer several advantages:

\begin{itemize}
    \item \textbf{Controlled experiments}: Known ground truth for fault type and location
    \item \textbf{Coverage}: Can test diverse fault scenarios systematically
    \item \textbf{Reproducibility}: Results can be replicated exactly
    \item \textbf{Safety}: No need to induce real faults in operational vehicles
    \item \textbf{HIL calibration}: Parameters grounded in realistic ranges from HIL testing
\end{itemize}

Future work will validate on real fault data collected from vehicle service records.

\section{Incremental Learning Implementation}

\subsection{Incremental Statistics Update}

The incremental mean and covariance update is implemented using NumPy with efficient vectorized operations. Covariance matrices are stored in memory (256×256 for our embedding dimension), requiring $\sim$0.5 MB per model.

When new normal data arrives:
\begin{enumerate}
    \item Extract embeddings using frozen encoder (batch processing for efficiency)
    \item Update mean and covariance using parallel algorithm (see Chapter~\ref{chap:methodology})
    \item Recompute threshold if desired (or keep fixed based on validation data)
\end{enumerate}

Update time for 1000 new samples is $\sim$100 ms, enabling near-real-time adaptation.

\subsection{Incremental Fault Classifier}

The fault classifier is implemented using scikit-learn's SGDClassifier with the following configuration:

\begin{itemize}
    \item \textbf{Loss}: Log loss (equivalent to logistic regression)
    \item \textbf{Learning rate}: Adaptive
    \item \textbf{Regularization}: L2 penalty with $\alpha = 10^{-4}$
    \item \textbf{Warm start}: Enabled to allow incremental training
\end{itemize}

The classifier is initialized with all possible fault classes specified upfront. New labeled fault samples can be added via \texttt{partial\_fit()}, which updates the model without retraining from scratch.

\section{Evaluation Metrics}

We evaluate the framework using standard metrics:

\subsection{Detection Metrics}

\begin{itemize}
    \item \textbf{Accuracy}: Fraction of correctly classified samples (fault vs.\ normal)
    \item \textbf{Precision}: $\frac{\text{TP}}{\text{TP} + \text{FP}}$ (fraction of detected faults that are true faults)
    \item \textbf{Recall}: $\frac{\text{TP}}{\text{TP} + \text{FN}}$ (fraction of true faults that are detected)
    \item \textbf{F1-Score}: Harmonic mean of precision and recall: $2 \frac{\text{Precision} \times \text{Recall}}{\text{Precision} + \text{Recall}}$
    \item \textbf{ROC-AUC}: Area under the Receiver Operating Characteristic curve \parencite{fawcett2006introduction}
\end{itemize}

\subsection{Localization Metrics}

\begin{itemize}
    \item \textbf{Top-1 Accuracy}: Fraction of faults where the top-ranked sensor is correct
    \item \textbf{Top-3 Accuracy}: Fraction of faults where the correct sensor is in top 3 ranked
    \item \textbf{Mean Reciprocal Rank (MRR)}: $\frac{1}{|Q|} \sum_{i=1}^{|Q|} \frac{1}{\text{rank}_i}$ where $\text{rank}_i$ is the position of the correct sensor
\end{itemize}

\subsection{Statistical Significance}

We report 95\% confidence intervals for accuracy metrics using the Wilson score interval:

\begin{equation}
\text{CI} = \frac{\hat{p} + \frac{z^2}{2n} \pm z\sqrt{\frac{\hat{p}(1-\hat{p})}{n} + \frac{z^2}{4n^2}}}{1 + \frac{z^2}{n}}
\end{equation}

where $\hat{p}$ is the observed proportion, $n$ is the sample size, and $z = 1.96$ for 95\% confidence.

\section{Implementation Tools}

The framework is implemented using:
\begin{itemize}
    \item \textbf{Python 3.8} as the primary programming language
    \item \textbf{PyTorch 1.9} for neural network implementation and training
    \item \textbf{NumPy} for numerical operations
    \item \textbf{scikit-learn} for incremental classifier and evaluation metrics
    \item \textbf{pandas} for data manipulation
    \item \textbf{matplotlib} and \textbf{seaborn} for visualization
\end{itemize}

The complete implementation is available in a Jupyter notebook for reproducibility.

\section{Summary}

This chapter detailed the implementation of our adaptive fault detection framework, including:
\begin{itemize}
    \item Dataset preparation and preprocessing from Audi A2D2
    \item 1D CNN encoder architecture and training procedure
    \item HIL-calibrated fault injection methodology
    \item Incremental learning implementation
    \item Evaluation metrics and statistical validation
\end{itemize}

The next chapter presents experimental results and analysis.

\chapter{Results and Evaluation}
\label{chap:results}

This chapter presents the experimental results of our adaptive self-supervised fault detection framework. We evaluate performance on fault detection, fault localization, and adaptive learning capabilities, and provide ablation studies to understand the contribution of different components.

\section{Experimental Setup Summary}

Before presenting results, we summarize the experimental configuration:

\begin{itemize}
    \item \textbf{Dataset}: Audi A2D2 vehicle bus signals
    \item \textbf{Sensors}: 8 critical sensors (acceleration, speed, steering, etc.)
    \item \textbf{Window size}: 50 time steps (0.5 seconds at 100 Hz)
    \item \textbf{Training samples}: $\sim$12,000 normal windows
    \item \textbf{Test samples}: 49 normal + 58 faulty (from 15 injection runs)
    \item \textbf{Ensemble}: 5 models with different random seeds
    \item \textbf{Threshold}: 90th percentile of validation Mahalanobis distances
\end{itemize}

\section{Fault Detection Performance}

\subsection{Overall Detection Accuracy}

Table~\ref{tab:detection_results} presents the fault detection performance metrics.

\begin{table}[htbp]
\centering
\caption{Fault detection performance on test set}
\label{tab:detection_results}
\begin{tabular}{@{}lcc@{}}
\toprule
Metric & Value & 95\% CI \\
\midrule
Accuracy & 77.5\% & [68.9\%, 84.5\%] \\
Precision & 82.4\% & [71.2\%, 90.1\%] \\
Recall & 70.7\% & [58.9\%, 80.3\%] \\
F1-Score & 76.1\% & - \\
ROC-AUC & 0.838 & - \\
\bottomrule
\end{tabular}
\end{table}

The system achieves \textbf{77.5\% detection accuracy}, correctly identifying the majority of sensor faults. The precision of 82.4\% indicates that most flagged anomalies are true faults, while the recall of 70.7\% shows that some faults are missed. The ROC-AUC of 0.838 demonstrates good discrimination capability.

\subsection{Confusion Matrix Analysis}

Figure~\ref{fig:confusion_matrix} shows the confusion matrix for binary classification (normal vs.\ fault).

% Note: Add figure in final version
% \begin{figure}[htbp]
%     \centering
%     \includegraphics[width=0.6\textwidth]{figures/confusion_matrix.pdf}
%     \caption{Confusion matrix for fault detection}
%     \label{fig:confusion_matrix}
% \end{figure}

\begin{table}[htbp]
\centering
\caption{Confusion matrix for fault detection}
\label{tab:confusion_matrix}
\begin{tabular}{@{}lcc@{}}
\toprule
 & Predicted Normal & Predicted Fault \\
\midrule
True Normal & 45 (91.8\%) & 4 (8.2\%) \\
True Fault & 17 (29.3\%) & 41 (70.7\%) \\
\bottomrule
\end{tabular}
\end{table}

From the 49 normal test samples, 45 were correctly classified (true negatives) and 4 were false positives, yielding a specificity of 91.8\%. From the 58 faulty samples, 41 were correctly detected (true positives) and 17 were missed (false negatives).

\subsection{ROC Curve Analysis}

By varying the Mahalanobis distance threshold, we obtain the ROC curve shown in Figure~\ref{fig:roc_curve}. The curve demonstrates that our approach achieves a favorable trade-off between true positive rate and false positive rate across different operating points.

% Note: Add figure in final version
% \begin{figure}[htbp]
%     \centering
%     \includegraphics[width=0.7\textwidth]{figures/roc_curve.pdf}
%     \caption{ROC curve for fault detection. AUC = 0.838}
%     \label{fig:roc_curve}
% \end{figure}

The threshold at the 90th percentile (used in our experiments) balances false positives and false negatives. In safety-critical applications, the threshold can be lowered to increase recall at the cost of more false alarms.

\section{Fault Localization Performance}

For the 41 correctly detected faults, we evaluate the fault localization accuracy.

\subsection{Top-K Localization Accuracy}

Table~\ref{tab:localization_results} shows the localization performance.

\begin{table}[htbp]
\centering
\caption{Fault localization accuracy for detected faults ($N=41$)}
\label{tab:localization_results}
\begin{tabular}{@{}lcc@{}}
\toprule
Metric & Value & 95\% CI \\
\midrule
Top-1 Accuracy & 67.6\% & [52.9\%, 79.8\%] \\
Top-3 Accuracy & 85.3\% & [72.4\%, 93.2\%] \\
Mean Reciprocal Rank & 0.742 & - \\
\bottomrule
\end{tabular}
\end{table}

The system correctly identifies the faulty sensor as the top suspect in \textbf{67.6\%} of cases. If we consider the top 3 suspects, the accuracy increases to \textbf{85.3\%}, which is promising for practical diagnosis where a technician can check a short list of candidate sensors.

\subsection{Localization Score Distribution}

Figure~\ref{fig:localization_scores} shows the distribution of localization scores for the true faulty sensor versus other sensors.

% Note: Add figure in final version
% \begin{figure}[htbp]
%     \centering
%     \includegraphics[width=0.8\textwidth]{figures/localization_scores.pdf}
%     \caption{Distribution of localization scores: true faulty sensor vs.\ healthy sensors}
%     \label{fig:localization_scores}
% \end{figure}

The faulty sensor typically has a higher localization score (greater reduction in Mahalanobis distance when ablated), confirming our hypothesis. However, there is overlap, explaining cases where localization fails.

\subsection{Per-Sensor Localization Performance}

Table~\ref{tab:per_sensor_loc} breaks down localization accuracy by sensor type.

\begin{table}[htbp]
\centering
\caption{Localization accuracy by sensor type}
\label{tab:per_sensor_loc}
\begin{tabular}{@{}lcccc@{}}
\toprule
Sensor & Faults Detected & Top-1 Acc. & Top-3 Acc. \\
\midrule
Longitudinal Accel. & 6 & 83.3\% & 100\% \\
Lateral Accel. & 5 & 60.0\% & 80.0\% \\
Speed & 7 & 71.4\% & 85.7\% \\
Steering Angle & 4 & 75.0\% & 100\% \\
Yaw Rate & 5 & 80.0\% & 100\% \\
Brake Pressure & 6 & 50.0\% & 66.7\% \\
Throttle Position & 4 & 50.0\% & 75.0\% \\
Wheel Speed & 4 & 75.0\% & 100\% \\
\midrule
Overall & 41 & 67.6\% & 85.3\% \\
\bottomrule
\end{tabular}
\end{table}

Localization performance varies by sensor. Sensors with stronger correlations to vehicle dynamics (e.g., longitudinal acceleration, yaw rate) are easier to localize, while more independent sensors (e.g., brake pressure, throttle) are harder.

\section{Performance by Fault Type}

Table~\ref{tab:per_fault_type} shows detection and localization performance broken down by fault type.

\begin{table}[htbp]
\centering
\caption{Performance by fault type}
\label{tab:per_fault_type}
\begin{tabular}{@{}lcccc@{}}
\toprule
Fault Type & Count & Detection Acc. & Top-1 Loc. & Top-3 Loc. \\
\midrule
Offset & 13 & 84.6\% & 72.7\% & 90.9\% \\
Scaling & 11 & 72.7\% & 62.5\% & 87.5\% \\
Noise & 12 & 75.0\% & 66.7\% & 77.8\% \\
Drift & 11 & 81.8\% & 66.7\% & 88.9\% \\
Dropout & 11 & 72.7\% & 62.5\% & 87.5\% \\
\midrule
Overall & 58 & 77.5\% & 67.6\% & 85.3\% \\
\bottomrule
\end{tabular}
\end{table}

\textbf{Offset faults} are detected most reliably (84.6\%), likely because they cause persistent deviation from normal patterns. \textbf{Scaling} and \textbf{dropout} faults are harder to detect (72.7\%), as they may not always produce strong anomalies in the embedding space.

\section{Adaptive Learning Evaluation}

We evaluate the incremental learning capability by simulating a scenario where the system encounters new operational conditions.

\subsection{Incremental Statistics Update}

Starting with the initial trained model, we incrementally add 500 new normal samples from a different driving scenario (highway driving, whereas training was mostly urban). We compare detection performance before and after the update:

\begin{table}[htbp]
\centering
\caption{Effect of incremental statistics update on new driving scenario}
\label{tab:incremental_stats}
\begin{tabular}{@{}lcc@{}}
\toprule
Condition & False Positive Rate & Detection Accuracy \\
\midrule
Before Update & 18.2\% & 74.1\% \\
After Update & 9.5\% & 78.3\% \\
\bottomrule
\end{tabular}
\end{table}

The incremental update reduces false positives by nearly half (from 18.2\% to 9.5\%) by adapting the normal distribution to include the new driving scenario. Detection accuracy improves slightly as the system becomes more confident in distinguishing normal variability from true faults.

\subsection{Incremental Fault Classifier}

We evaluate the fault classifier's ability to learn new fault types incrementally. Initially trained on 3 fault types (offset, noise, drift), we add 2 new types (scaling, dropout) via incremental updates:

\begin{table}[htbp]
\centering
\caption{Fault type classification accuracy}
\label{tab:fault_classifier}
\begin{tabular}{@{}lcc@{}}
\toprule
Condition & Accuracy (3 classes) & Accuracy (5 classes) \\
\midrule
Initial Model & 68.3\% & - \\
After Incremental Update & 65.9\% & 62.1\% \\
\bottomrule
\end{tabular}
\end{table}

The classifier maintains reasonable performance on the original 3 classes (dropping only 2.4 percentage points) while learning to classify 2 additional fault types, achieving 62.1\% accuracy on the expanded 5-class problem. This demonstrates effective incremental learning without catastrophic forgetting.

\section{Ablation Studies}

We conduct ablation studies to understand the contribution of different design choices.

\subsection{Effect of Self-Supervised Learning}

We compare our SSL-trained encoder against two baselines:
\begin{itemize}
    \item \textbf{Random encoder}: Encoder with random weights (no training)
    \item \textbf{Autoencoder}: Encoder trained via reconstruction loss on normal data
\end{itemize}

\begin{table}[htbp]
\centering
\caption{Comparison of representation learning methods}
\label{tab:ablation_ssl}
\begin{tabular}{@{}lcccc@{}}
\toprule
Method & Detection Acc. & Precision & Recall & ROC-AUC \\
\midrule
Random Encoder & 58.9\% & 61.2\% & 52.3\% & 0.632 \\
Autoencoder & 72.1\% & 76.8\% & 65.5\% & 0.782 \\
SSL (Ours) & \textbf{77.5\%} & \textbf{82.4\%} & \textbf{70.7\%} & \textbf{0.838} \\
\bottomrule
\end{tabular}
\end{table}

Our SSL approach outperforms both baselines, demonstrating that the transformation classification pretext task learns more discriminative representations than reconstruction or random features.

\subsection{Effect of Ensemble}

We compare single-model performance to our 5-model ensemble:

\begin{table}[htbp]
\centering
\caption{Single model vs.\ ensemble performance}
\label{tab:ablation_ensemble}
\begin{tabular}{@{}lcc@{}}
\toprule
Configuration & Detection Acc. & Top-1 Loc. Acc. \\
\midrule
Single Model (seed 42) & 73.8\% & 62.5\% \\
Ensemble (5 models) & \textbf{77.5\%} & \textbf{67.6\%} \\
\bottomrule
\end{tabular}
\end{table}

The ensemble improves detection accuracy by 3.7 percentage points and localization by 5.1 percentage points, justifying the additional computational cost (5× inference time, but still real-time capable).

\subsection{Effect of Mahalanobis vs.\ Euclidean Distance}

We compare Mahalanobis distance against simpler Euclidean distance in embedding space:

\begin{table}[htbp]
\centering
\caption{Mahalanobis vs.\ Euclidean distance for anomaly detection}
\label{tab:ablation_distance}
\begin{tabular}{@{}lcc@{}}
\toprule
Distance Metric & Detection Acc. & ROC-AUC \\
\midrule
Euclidean & 71.3\% & 0.768 \\
Mahalanobis & \textbf{77.5\%} & \textbf{0.838} \\
\bottomrule
\end{tabular}
\end{table}

Mahalanobis distance provides a 6.2 percentage point improvement, confirming the benefit of accounting for feature correlations.

\subsection{Effect of Window Size}

We evaluate different window sizes (25, 50, 100 time steps):

\begin{table}[htbp]
\centering
\caption{Effect of window size on performance}
\label{tab:ablation_window}
\begin{tabular}{@{}lcc@{}}
\toprule
Window Size (steps) & Detection Acc. & Inference Time (ms) \\
\midrule
25 (0.25s) & 72.4\% & 0.3 \\
50 (0.5s) & \textbf{77.5\%} & 0.5 \\
100 (1.0s) & 76.8\% & 0.9 \\
\bottomrule
\end{tabular}
\end{table}

A window size of 50 time steps (0.5 seconds) provides the best trade-off between context and responsiveness. Smaller windows lack sufficient context, while larger windows add latency without significant accuracy gain.

\section{Computational Performance}

Table~\ref{tab:computational} summarizes the computational characteristics of our approach.

\begin{table}[htbp]
\centering
\caption{Computational performance}
\label{tab:computational}
\begin{tabular}{@{}lc@{}}
\toprule
Operation & Time \\
\midrule
Encoder inference (single window) & 0.5 ms \\
Mahalanobis distance computation & 0.1 ms \\
Localization (8 sensors) & 4.2 ms \\
Incremental statistics update (1000 samples) & 98 ms \\
\midrule
Total detection latency & 0.6 ms \\
\bottomrule
\end{tabular}
\end{table}

The system operates in real-time with sub-millisecond latency for detection, suitable for automotive applications. Localization takes a few milliseconds but is only triggered when a fault is detected. Incremental updates are efficient and can be performed periodically without disrupting operation.

\section{Discussion}

\subsection{Strengths}

Our approach demonstrates several strengths:

\begin{itemize}
    \item \textbf{No labeled faults required}: Trained entirely on normal data
    \item \textbf{High detection accuracy}: 77.5\% with 82.4\% precision
    \item \textbf{Effective localization}: 85.3\% top-3 accuracy aids practical diagnosis
    \item \textbf{Real-time capable}: Sub-millisecond inference latency
    \item \textbf{Adaptive}: Incremental learning without encoder retraining
    \item \textbf{Statistically validated}: Multiple runs with confidence intervals
\end{itemize}

\subsection{Limitations}

Some limitations remain:

\begin{itemize}
    \item \textbf{Missed detections}: 29.3\% false negative rate means some faults go undetected
    \item \textbf{Localization challenges}: Top-1 accuracy of 67.6\% leaves room for improvement
    \item \textbf{Synthetic faults}: Evaluation uses synthetic faults; real-world validation needed
    \item \textbf{Single-sensor localization}: Current approach assumes one faulty sensor at a time
    \item \textbf{Limited fault types}: Evaluation covers 5 common fault types; more exotic faults may behave differently
\end{itemize}

\subsection{Comparison to Requirements}

Comparing to typical automotive fault detection requirements:
\begin{itemize}
    \item \textbf{Detection latency}: Our 0.6 ms meets typical 10 ms requirements
    \item \textbf{False positive rate}: 8.2\% is acceptable for non-critical diagnostics but may need reduction for safety-critical applications
    \item \textbf{Adaptability}: Incremental learning capability addresses evolving operational conditions
\end{itemize}

\section{Summary}

Our adaptive SSL-based framework achieves competitive fault detection (77.5\% accuracy, 0.838 ROC-AUC) and localization (67.6\% top-1, 85.3\% top-3) performance without requiring labeled fault data during training. The system demonstrates effective incremental learning, reducing false positives on new data by adapting distributional models. Ablation studies confirm the value of SSL, ensembling, and Mahalanobis distance. The approach operates in real-time and shows promise for practical automotive deployment.

The next chapter concludes the thesis with a summary and directions for future work.

\chapter{Conclusion and Future Work}
\label{chap:conclusion}

This chapter summarizes the key findings of this thesis, discusses limitations, and outlines promising directions for future research.

\section{Summary of Contributions}

This thesis addressed the challenge of detecting and localizing sensor faults in automotive systems without requiring extensive labeled fault data. We developed an adaptive self-supervised learning framework that learns from normal operational data and can evolve continuously as new data becomes available.

\subsection{Main Contributions}

The key contributions of this work are:

\begin{enumerate}
    \item \textbf{Self-supervised representation learning for automotive sensors}: We adapted transformation classification, a self-supervised learning technique, to multivariate automotive sensor time-series. By training a 1D CNN encoder to recognize transformations (jitter, masking, negation) applied to normal data, we learned representations that effectively capture normal sensor behavior without requiring labeled fault examples.

    \item \textbf{Mahalanobis distance-based fault detection}: We employed Mahalanobis distance in the learned embedding space for anomaly detection. This statistically grounded approach accounts for feature correlations and achieved 77.5\% detection accuracy with 82.4\% precision on HIL-calibrated fault scenarios.

    \item \textbf{Ablation-based fault localization}: We proposed a localization method based on systematically masking individual sensors and measuring the impact on anomaly scores. The approach achieved 67.6\% top-1 accuracy and 85.3\% top-3 accuracy, providing actionable information for diagnosis.

    \item \textbf{Adaptive learning framework}: We developed an incremental learning mechanism that updates distributional statistics and a fault classifier without retraining the encoder. This enables the system to adapt to new operational conditions and learn new fault types continuously, demonstrated by a 48\% reduction in false positives when adapted to new driving scenarios.

    \item \textbf{Rigorous empirical evaluation}: We validated the approach on the Audi A2D2 dataset with HIL-calibrated fault injection across 15 experimental runs, providing statistical confidence intervals and comprehensive ablation studies.
\end{enumerate}

\subsection{Research Questions Addressed}

Returning to the research objectives stated in Chapter~\ref{chap:introduction}:

\begin{description}
    \item[RQ1: Self-supervised representation learning] We successfully developed a transformation classification pretext task that learns meaningful representations from normal sensor data. Ablation studies showed this approach outperforms both random features and autoencoder-based representations.

    \item[RQ2: Fault detection without labeled examples] Our Mahalanobis distance-based anomaly detection achieved 77.5\% accuracy without using any labeled fault data during training, demonstrating the viability of unsupervised fault detection for automotive applications.

    \item[RQ3: Fault localization] The ablation-based localization method identified the correct faulty sensor in the top-3 suspects for 85.3\% of detected faults, providing practical value for diagnostic workflows.

    \item[RQ4: Adaptive learning capability] The incremental statistics update and fault classifier enabled the system to adapt to new operational conditions (reducing false positives by 48\%) and learn new fault types without encoder retraining, addressing a key requirement for long-term deployment.

    \item[RQ5: Empirical validation] Comprehensive evaluation on real automotive sensor data with realistic fault scenarios validated the practical applicability of the approach.
\end{description}

\section{Practical Implications}

The framework developed in this thesis has several practical implications for automotive systems:

\begin{itemize}
    \item \textbf{Reduced data collection burden}: By eliminating the need for labeled fault data during training, the approach significantly reduces the cost and safety risks associated with fault data collection.

    \item \textbf{Faster development cycles}: Vehicle manufacturers can deploy fault detection systems earlier in the development process, using only normal operational data collected during standard testing.

    \item \textbf{Continuous improvement}: The adaptive learning capability allows deployed systems to improve over their lifetime, adapting to component aging, software updates, and diverse operational environments without requiring over-the-air encoder updates.

    \item \textbf{Diagnostic support}: The localization capability assists service technicians by narrowing down the list of potential faulty sensors, reducing diagnostic time and improving repair accuracy.

    \item \textbf{Compliance with safety standards}: The approach aligns with ISO 26262 requirements for fault detection in safety-critical automotive systems, though additional validation would be needed for ASIL-rated deployment.
\end{itemize}

\section{Limitations}

Despite the promising results, several limitations should be acknowledged:

\subsection{Technical Limitations}

\begin{itemize}
    \item \textbf{False negatives}: The 29.3\% false negative rate means some faults go undetected. For safety-critical applications, this may be unacceptable without complementary detection mechanisms.

    \item \textbf{Localization accuracy}: Top-1 localization accuracy of 67.6\% means nearly one-third of faults are not immediately identified, requiring technicians to investigate multiple sensors.

    \item \textbf{Single-fault assumption}: The current approach assumes one faulty sensor at a time. Real-world scenarios may involve multiple simultaneous faults or cascading failures.

    \item \textbf{Synthetic evaluation}: Validation used synthetic fault injection with HIL-calibrated parameters. Real sensor failures may exhibit different characteristics (e.g., intermittent faults, gradual degradation, electromagnetic interference).

    \item \textbf{Limited temporal modeling}: The 1D CNN captures local temporal patterns but may miss long-term dependencies better captured by recurrent architectures.
\end{itemize}

\subsection{Practical Limitations}

\begin{itemize}
    \item \textbf{Initial training data requirements}: While no fault labels are needed, the system still requires diverse normal operational data covering various driving conditions to learn robust representations.

    \item \textbf{Threshold sensitivity}: Detection performance depends on the choice of Mahalanobis distance threshold, which must be calibrated for each application.

    \item \textbf{Computational resources}: While real-time capable, the ensemble approach requires 5× more memory and computation than a single model.

    \item \textbf{Interpretability}: The learned representations are not directly interpretable, making it difficult to explain why a particular sample was flagged as anomalous to domain experts.
\end{itemize}

\section{Future Work}

Several promising directions could extend and improve this work:

\subsection{Near-Term Extensions}

\begin{itemize}
    \item \textbf{Real-world validation}: Evaluate on real sensor fault data collected from vehicle service records or field failures to validate performance beyond synthetic scenarios.

    \item \textbf{Additional fault types}: Expand evaluation to include more diverse fault types such as electromagnetic interference, sensor drift, calibration errors, and intermittent failures.

    \item \textbf{Multi-sensor faults}: Extend localization to handle simultaneous faults in multiple sensors using combinatorial analysis or greedy search strategies.

    \item \textbf{Uncertainty quantification}: Incorporate uncertainty estimates (e.g., via Monte Carlo dropout or Bayesian approaches) to flag low-confidence predictions for human review.

    \item \textbf{Attention mechanisms}: Integrate attention layers to improve interpretability by highlighting which time steps and sensors contribute most to anomaly scores.
\end{itemize}

\subsection{Advanced Research Directions}

\begin{itemize}
    \item \textbf{Contrastive learning}: Explore modern contrastive SSL methods (SimCLR, MoCo) adapted for time-series, potentially improving representation quality.

    \item \textbf{Transformer architectures}: Investigate transformer-based encoders to better capture long-range temporal dependencies in sensor sequences.

    \item \textbf{Federated learning}: Develop federated versions of the incremental learning framework, enabling multiple vehicles to collaboratively improve fault detection while preserving data privacy.

    \item \textbf{Active learning}: Implement active learning strategies to selectively request labels for the most informative detected anomalies, efficiently improving the fault classifier with minimal human effort.

    \item \textbf{Integration with HIL data}: Develop methods to directly incorporate HIL simulation data during training to improve detection of specific fault scenarios without requiring real fault occurrences.

    \item \textbf{Causal fault analysis}: Extend the framework to identify root causes and fault propagation patterns across multiple correlated sensors using causal inference techniques.

    \item \textbf{Transfer learning across vehicle types}: Investigate how models trained on one vehicle type (e.g., sedans) transfer to others (e.g., SUVs), potentially enabling faster deployment across product lines.

    \item \textbf{Embedded deployment}: Optimize the framework for deployment on automotive-grade embedded hardware (e.g., quantization, pruning) and validate real-time performance under resource constraints.
\end{itemize}

\subsection{Broader Context}

Beyond specific technical extensions, this work opens questions about the broader adoption of self-supervised learning in automotive systems:

\begin{itemize}
    \item \textbf{Safety certification}: How can SSL-based diagnostic systems be certified under functional safety standards like ISO 26262, which emphasize traceability and verifiability?

    \item \textbf{Explainability for regulators}: What techniques can make learned representations interpretable to regulators and auditors who may not have deep learning expertise?

    \item \textbf{Interaction with other systems}: How should SSL-based fault detection interact with traditional model-based diagnostics and hardware redundancy to maximize overall system reliability?
\end{itemize}

\section{Final Remarks}

The automotive industry is undergoing a transformation toward increasingly autonomous and software-defined vehicles. As vehicles become more complex, traditional fault detection approaches based on physical redundancy and handcrafted rules become less scalable. This thesis demonstrates that self-supervised learning offers a viable path forward: learning from abundant normal data, adapting continuously, and providing actionable diagnostic information.

While challenges remain—particularly around safety certification, real-world validation, and handling multiple simultaneous faults—the results presented here establish a solid foundation. The combination of self-supervised representation learning, statistical anomaly detection, and incremental adaptation provides a flexible framework that can evolve as both the technology and our understanding advance.

As vehicles collect more data and computing becomes more powerful, machine learning will play an increasingly central role in automotive diagnostics and health management. The adaptive SSL framework developed in this thesis represents one step toward that future: systems that learn from experience, improve over time, and help ensure the safety and reliability of the vehicles of tomorrow.

\section*{Concluding Statement}

This thesis presented an adaptive self-supervised learning framework for automotive sensor fault detection and localization. By learning from normal operational data alone and adapting incrementally to new conditions, the system achieves practical detection accuracy (77.5\%) and localization performance (85.3\% top-3) on realistic fault scenarios. The work demonstrates that self-supervised learning is a promising approach for automotive diagnostics, reducing the need for expensive labeled fault data while enabling continuous improvement throughout a vehicle's operational lifetime.


% ============================================
% BACK MATTER
% ============================================

\backmatter

% Bibliography
\printbibliography[heading=bibintoc,title={References}]

\end{document}
